\begin{theorem}[容量亏损指数的 Mellin--Laplace 变分与冻结暴露平面]\label{thm:xi-capacity-defect-exponent-mellin-laplace-face}
沿用定理 \ref{thm:xi-dyadic-oracle-capacity-mellin-recovery} 的记号，并设
\[
\alpha_*:=\lim_{m\to\infty}\frac{1}{m}\log M_m,
\qquad
g_*:=\lim_{m\to\infty}\frac{1}{m}\log K_m
\]
存在。对每个 \(\beta\in[0,\alpha_*]\)，定义容量亏损指数近似
\[
\Lambda_m(\beta):=
\frac{1}{m}\log\Bigl(2^m-\mathcal C_m^{\mathrm{cont}}(e^{\beta m})\Bigr).
\]
进一步假设极限
\[
\Lambda(\beta):=\lim_{m\to\infty}\Lambda_m(\beta)
\]
对每个 \(\beta\in[0,\alpha_*]\) 存在，且 \(\Lambda_m\to\Lambda\) 在 \([0,\alpha_*]\) 的每个紧子区间上一致收敛。则对任意实数 \(a>1\)，极限
\[
P(a):=\lim_{m\to\infty}\frac{1}{m}\log S_a(m)
\]
存在，并满足精确变分公式
\[
\boxed{\;
P(a)=\sup_{0\le \beta\le \alpha_*}\Bigl((a-1)\beta+\Lambda(\beta)\Bigr).\;}
\]
若进一步设
\[
\alpha_*^{(2)}:=\limsup_{m\to\infty}\frac{1}{m}\log M_m^{(2)}<\alpha_*,
\]
则对每个
\[
\alpha_*^{(2)}<\beta<\alpha_*
\]
都成立
\[
\boxed{\;
\Lambda(\beta)=\alpha_*+g_*.\;}
\]
因此容量亏损指数在严格谱隙区间上形成常值暴露平面；并且当
\[
a>a_{\mathrm c}:=\frac{\log\varphi-g_*}{\alpha_*-\alpha_*^{(2)}}
\]
时，冻结线性支
\[
P(a)=a\alpha_*+g_*
\]
正由该常值平面的右端在 Laplace 极限中暴露出来。
\end{theorem}
\begin{proof}
由定理 \ref{thm:xi-dyadic-oracle-capacity-mellin-recovery}，对任意 \(a>1\) 都有
\[
S_a(m)
=
a(a-1)\int_0^\infty T^{a-2}\Bigl(2^m-\mathcal C_m^{\mathrm{cont}}(T)\Bigr)\,\dd T.
\]
又因
\[
\mathcal C_m^{\mathrm{cont}}(T)=2^m
\qquad (T\ge M_m),
\]
故积分可截断到 \([0,M_m]\)。令 \(T=e^{\beta m}\)，则 \(\dd T=me^{\beta m}\dd\beta\)，从而
\[
S_a(m)
=
a(a-1)m
\int_0^{\frac1m\log M_m}
\exp\!\Bigl(
m\bigl((a-1)\beta+\Lambda_m(\beta)\bigr)
\Bigr)\,\dd\beta.
\]
由于
\[
\frac1m\log M_m\to \alpha_*,
\]
由于 \(\Lambda_m\to\Lambda\) 在 \([0,\alpha_*]\) 上可由紧性统一控制，故可直接对上式的指数积分应用 Laplace 原理，得到
\[
\lim_{m\to\infty}\frac1m\log S_a(m)
=
\sup_{0\le \beta\le \alpha_*}
\Bigl((a-1)\beta+\Lambda(\beta)\Bigr).
\]
这就得到了所示变分公式。

再设 \(\alpha_*^{(2)}<\beta<\alpha_*\)。由定义，
\[
\frac1m\log M_m^{(2)}\le \alpha_*^{(2)}+o(1),
\qquad
\frac1m\log M_m=\alpha_*+o(1),
\]
故对充分大的 \(m\) 必有
\[
M_m^{(2)}<e^{\beta m}<M_m.
\]
于是
\[
2^m-\mathcal C_m^{\mathrm{cont}}(e^{\beta m})
=
\sum_{x\in X_m}\bigl(d_m(x)-e^{\beta m}\bigr)_+
=
K_m\bigl(M_m-e^{\beta m}\bigr),
\]
因为只有最大纤维对上式仍有贡献。又由 \(\beta<\alpha_*\) 得
\[
\frac{e^{\beta m}}{M_m}
=
\exp\!\bigl(-(\alpha_*-\beta)m+o(m)\bigr)\to 0,
\]
故
\[
M_m-e^{\beta m}=M_m(1+o(1)).
\]
从而
\[
\Lambda(\beta)
=
\lim_{m\to\infty}\frac1m\log\bigl(K_mM_m(1+o(1))\bigr)
=
\alpha_*+g_*.
\]

最后，在上述常值区间上
\[
(a-1)\beta+\Lambda(\beta)
=
(a-1)\beta+\alpha_*+g_*,
\]
其值随 \(\beta\) 严格递增。因而该表达式在
\(
\beta\uparrow \alpha_*
\)
时趋于
\[
a\alpha_*+g_*.
\]
这与定理 \ref{thm:fold-pressure-freezing-threshold} 在 \(a>a_{\mathrm c}\) 时给出的冻结公式一致，故所述线性支正由该常值暴露平面的右端在 Laplace 极限中暴露出来。证毕。
\end{proof}

\endinput
